\section{The Bailout Protocol}\label{sec:bailout}

The Bailout Protocol is the mandatory exception-escalation mechanism of the \bxthree{} Framework. It is not an error-handling routine selected at runtime, nor an optional escalation pathway available to a system that chooses to invoke it. It is an architectural compulsion: any unresolved exception state encountered by any node in a \bxthree{} system \emph{must} propagate upward through the accountability chain until it reaches a named human principal, and no machine actor along that chain may absorb, resolve, or suppress the exception. This section provides the complete formal specification: the deterministic state machine governing exception propagation, the precise conditions governing each transition, the bail-out trigger condition, the escalation algorithm, the bypass mechanism enforcing the no-machine-actor property, and the timeout handling that guarantees liveness under all failure conditions.

The protocol operates against the backdrop of the \bxthree{} three-layer model. The \purpose{} layer carries intent, judgment, and final authority; it defines the operating range $R(n)$ for each node $n$ and receives all escalated exceptions. The \bounds{} layer carries bounded reasoning and constrained execution within $R(n)$; it evaluates whether proposed actions fall within the node's provisioned authority and initiates bailout when they do not. The \fact{} layer carries deterministic enforcement and forensic auditability; it verifies proposed actions against the authorization ledger, enforces the state machine's transition relation, and maintains the append-only Forensic Ledger as an immutable record of every protocol event. Each node carries an immutable parent pointer $\alpha(n)$ established at provisioning and a reference to its human accountability anchor $h(n)$ — a named human principal — that is similarly immutable for the node's operational lifetime. The Bailout Protocol exploits this structure to guarantee that exception states propagate to $h(n)$ without passing through any intermediate machine actor as a terminus.

% --------------------------------------------------------------
\subsection{State Machine Definition}\label{sec:state-machine}
% --------------------------------------------------------------

The Bailout Protocol is formally modeled as a deterministic finite state machine embedded in the \fact{} layer of every \bxthree{} node:

\[
\mathcal{B} \;=\; (Q,\; q_0,\; \Sigma,\; \rightarrow,\; Q_F)
\]

\begin{itemize}\itemsep=0pt
\item[$Q$] $\{\; \text{IDLE},\; \text{BOUNDED},\; \text{BAIL\_PENDING},\; \text{ESCALATING},\; \text{HUMAN\_ACKNOWLEDGED},\; \text{RESOLVED}\; \}$ is the finite set of protocol states.
\item[$q_0$] $\text{IDLE}$ is the designated initial state, entered at node startup and after every successful directive resolution.
\item[$\Sigma$] The input alphabet comprising trigger events $e_{\text{trigger}}$, authorization confirmations $e_{\text{auth}}$, timeout signals $e_{\text{timeout}}$, human acknowledgments $e_{\text{ack}}$, binding resolution directives $e_{\text{resolve}}$, and rejection directives $e_{\text{reject}}$.
\item[$\rightarrow$] $\; \subseteq Q \times \Sigma \times Q$ is the deterministic transition relation (Table~\ref{tab:bailout-transition}).
\item[$Q_F$] $\{\text{RESOLVED}\}$ is the set of terminal accepting states.
\end{itemize}

A node's state $q \in Q$ is stored in the node's \fact{} layer worksheet and is observable by its parent node and by the Bailout Monitor. State updates are atomic: no intermediate or partially-executing states are observable externally. The machine is \emph{deterministic}: for every $(q, \sigma) \in Q \times \Sigma$, at most one $q' \in Q$ satisfies $(q, \sigma, q') \in \rightarrow$. Determinism is enforced by the Bailout Monitor's priority function $\pi$, which totally orders simultaneous trigger conditions and selects the highest-priority enabled transition before committing any state change. If no transition is defined for a $(q, \sigma)$ pair, the machine stalls and the Bailout Monitor records a protocol violation in the Forensic Ledger.

\begin{table}[htbp]
\caption{Bailout Protocol state transition relation $\rightarrow \subseteq Q \times \Sigma \times Q$. Rows are current state; columns are input events. Entries denote the next state. Blank cells indicate that no transition is defined — the machine stalls and the Bailout Monitor records a protocol violation.}
\label{tab:bailout-transition}
\begin{tabular}{@{}lcccccc@{}}
\toprule
 & \multicolumn{6}{c}{Input event} \\
\cmidrule(l){2-7}
Current state & $e_{\text{trigger}}$ & $e_{\text{auth}}$ & $e_{\text{timeout}}$ & $e_{\text{ack}}$ & $e_{\text{resolve}}$ & $e_{\text{reject}}$ \\
\midrule
IDLE              & BOUNDED            & — & — & — & — & — \\
BOUNDED           & — & HUMAN\_ACKNOWLEDGED & BAIL\_PENDING & — & — & — \\
BAIL\_PENDING     & — & — & ESCALATING & — & — & — \\
ESCALATING        & — & HUMAN\_ACKNOWLEDGED & ESCALATING & IDLE & RESOLVED & RESOLVED \\
HUMAN\_ACKNOWLEDGED & — & — & — & IDLE & RESOLVED & IDLE \\
RESOLVED          & — & — & — & — & — & — \\
\bottomrule
\end{tabular}
\end{table}

\bigskip
\noindent\textbf{State definitions.}

\begin{itemize}\itemsep=0pt
\item[\textbf{IDLE.}] The node is operating within its provisioned operating range $R(n)$. No exception condition is active. The \fact{} layer dispatches normal action requests through the standard execution path. The node may transition to BOUNDED upon receipt of a trigger event satisfying the bail-out trigger condition (TC), defined in \S~\ref{sec:bailout-trigger}.

\item[\textbf{BOUNDED.}] An exception $e$ has been detected by the \bounds{} layer and is being held at the \bounds{}--\fact{} layer boundary. A time-bounded resolution attempt is in progress. All resolution paths within current authority remain open. If the Bounds Engine produces a resolution $r$ that the \fact{} layer approves before $T_{\text{bounds}}$ expires, the node returns to IDLE. If all resolution paths are exhausted or $T_{\text{bounds}}$ expires, the state advances to BAIL\_PENDING.

\item[\textbf{BAIL\_PENDING.}] The node has determined — either through an explicit \texttt{bailout\_signal} or through $T_{\text{bounds}}$ expiration without a \fact-layer-approved resolution — that the exception exceeds its current authority. The Bailout Protocol is formally activated. This is the last architectural point at which a machine actor could theoretically suppress escalation. It cannot: entry to BAIL\_PENDING triggers an immediate atomic transition to ESCALATING via $e_{\text{timeout}}$ (T4), with no dwell time and no machine-actor intervention permitted.

\item[\textbf{ESCALATING.}] The exception is propagating upward toward $h(n)$. Each intermediate node receives the exception, appends its diagnostic context to the propagation chain, and forwards the exception via \texttt{SendToParent} without attempting local resolution. The state remains ESCALATING until either (a) the human anchor acknowledges receipt ($e_{\text{ack}}$: T5), (b) the human anchor issues a binding directive ($e_{\text{resolve}}$ or $e_{\text{reject}}$: T7), or (c) all retry pathways timeout after $N_{\max}$ retries at $T_{\text{cap}}$.

\item[\textbf{HUMAN\_ACKNOWLEDGED.}] The human anchor has received the exception, observed the diagnostic context, and acknowledged it. The node awaits a directive. No machine actor may issue, simulate, or substitute a directive in this state. The node remains until the human principal issues \texttt{ACK} (return to IDLE) or a binding directive \texttt{RESOLVE} / \texttt{REJECT} (enter RESOLVED).

\item[\textbf{RESOLVED.}] The human principal has issued a directive and the protocol has recorded it in the authorization ledger. The directive is one of: \texttt{ACK} (continue with current parameters), \texttt{RESOLVE} (execute a specific binding resolution), or \texttt{REJECT} (disallow the proposed action and enter a quiescent error state). This state is terminal for the current protocol run.
\end{itemize}

\bigskip
\noindent\textbf{State invariants.} The following invariants hold for all reachable states of $\mathcal{B}$:

\begin{align}
\text{INV}_1 &: \forall n \in \text{nodes},\; \text{state}(n) \in Q \setminus \{\text{ESCALATING}\} \iff \neg\text{active\_bailout}(n). \tag{INV-1}
\end{align}

That is, a node occupies a non-escalating state \emph{iff} it has no active bailout propagation in flight. A node may carry a bailout history (a ledger entry recording a past propagation) while in IDLE, but no active propagation may be underway.

\begin{align}
\text{INV}_2 &: \text{state}(n) = \text{BAIL\_PENDING} \implies \text{active\_bailout}(n) \land \text{transition}(\text{BAIL\_PENDING}, e_{\text{timeout}}) = \text{ESCALATING}. \tag{INV-2}
\end{align}

That is, any node in BAIL\_PENDING has an active bailout and the only permissible next transition is to ESCALATING, enforced by the Bailout Monitor without machine-actor discretion.

% --------------------------------------------------------------
\subsection{Transition Conditions}\label{sec:transitions}
% --------------------------------------------------------------

Each transition T1–T7 encodes a precise condition on the protocol state and the \bxthree{} layer interfaces:

\begin{itemize}\itemsep=0pt
\item[T1 --- \texttt{IDLE $\rightarrow$ BOUNDED}:] Fires when $\text{TRIGGER}(n, x)$ holds for some $x$ in the node's observable input space. This is the entry point for all bailout activations and is the sole permitted entry; no transition skips directly to ESCALATING or HUMAN\_ACKNOWLEDGED without passing through BOUNDED.

\item[T2 --- \texttt{BOUNDED $\rightarrow$ HUMAN\_ACKNOWLEDGED}:] Fires when the \bounds{} layer produces a resolution $r$ that the \fact{} layer approves via the authorization ledger before $T_{\text{bounds}}$ expires. Approval requires that $r \in R(n)$ and that $r$ is not explicitly flagged by the Bounds Engine as requiring human review. This is the \emph{success} path: the node handles the exception autonomously and returns to normal operation.

\item[T3 --- \texttt{BOUNDED $\rightarrow$ BAIL\_PENDING}:] Fires when either (a) $T_{\text{bounds}}$ expires without a \fact-layer-approved resolution, or (b) the Bounds Engine emits an explicit \texttt{bailout\_signal}. This transition marks the formal activation of the Bailout Protocol at the originating node.

\item[T4 --- \texttt{BAIL\_PENDING $\rightarrow$ ESCALATING}:] Fires atomically upon entry to BAIL\_PENDING via the $e_{\text{timeout}}$ signal generated by the \fact{} layer timer. No machine-actor dwell time is permitted: the transition is immediate, enforced by the Bailout Monitor pre-commit hook. This is the architectural no-suppression gate.

\item[T5 --- \texttt{ESCALATING $\rightarrow$ HUMAN\_ACKNOWLEDGED}:] Fires when $h(n)$ acknowledges receipt of the exception via the \texttt{ACK} directive. Acknowledgment confirms that the human has observed the exception and is actively engaged.

\item[T6 --- \texttt{ESCALATING $\rightarrow$ IDLE}:] Fires when $h(n)$ issues a \texttt{REJECT} directive during propagation. The exception is logged and discarded; no resolution is recorded. The node returns to IDLE with a quiescent error tag in the Forensic Ledger.

\item[T7 --- \texttt{ESCALATING $\rightarrow$ RESOLVED}:] Fires when $h(n)$ issues either \texttt{RESOLVE} or \texttt{REJECT} after acknowledgment. The directive is recorded in the authorization ledger and the protocol run terminates.

\item[T8 --- \texttt{HUMAN\_ACKNOWLEDGED $\rightarrow$ IDLE}:] Fires when $h(n)$ issues \texttt{ACK} after acknowledgment, returning the node to normal operation without a binding resolution.

\item[T9 --- \texttt{HUMAN\_ACKNOWLEDGED $\rightarrow$ RESOLVED}:] Fires when $h(n)$ issues \texttt{RESOLVE} or \texttt{REJECT}, terminating the protocol with a binding directive recorded in the ledger.
\end{itemize}

All transitions are recorded in the Forensic Ledger with the node identifier, timestamp, triggering event, and source layer. The ledger entry is immutable and cryptographically sealed upon write.

% --------------------------------------------------------------
\subsection{Bail-Out Trigger Condition}\label{sec:bailout-trigger}
% --------------------------------------------------------------

The trigger condition TC is the mechanism by which a node enters BOUNDED state (T1). Rather than requiring system designers to enumerate specific error codes in advance — a practice that provides coverage only for anticipated failure modes and leaves unanticipated ones invisible — the Bailout Protocol defines the trigger in terms of \emph{observational boundary violation}: a node escalates precisely when its observable input space contains a condition outside its provisioned operating range $R(n)$ and the node cannot produce a \fact-layer-approved resolution within $T_{\text{bounds}}$.

\bigskip
\noindent\textbf{Primary trigger condition.} The trigger $\text{TRIGGER}(n, x)$ is defined as a conjunction of three sub-conditions, all of which must hold simultaneously:

\begin{equation}
\text{TRIGGER}(n, x) \;\equiv\; \underbrace{x \notin R(n)}_{(i)} \;\land\; \underbrace{\neg\text{resolved}(n, x)}_{(ii)} \;\land\; \underbrace{\text{confidence}(n, x) < \tau}_{(iii)} \tag{TC}
\end{equation}

where:

\begin{itemize}\itemsep=0pt
\item[(i)] $x$ is an element of the node's observable input space that falls outside the provisioned operating range $R(n)$. This is the \emph{bound signal}: the node observes something it was not provisioned to handle autonomously.

\item[(ii)] $x$ has not been previously resolved by the node within the current operational session. This prevents re-trigger on known-exception inputs that have already received human adjudication.

\item[(iii)] $\text{confidence}(n, x) < \tau$ is the node's self-assessed confidence that it can produce a correct resolution for $x$ without human input, where $\tau$ is the provisioned confidence threshold. The confidence function $\text{confidence} : \text{node} \times \text{input} \rightarrow [0, 1]$ is defined at node provisioning and is \emph{not} adjustable at runtime by any machine actor. This prevents a machine actor from inflating its own confidence to suppress a warranted escalation.
\end{itemize}

\bigskip
\noindent\textbf{Trigger priority function.} When multiple trigger conditions fire simultaneously for node $n$, they are totally ordered by the priority function $\pi(n, x)$:

\begin{equation}
\pi(n, x) \;=\; w_1 \cdot \text{severity}(x) \;+\; w_2 \cdot \bigl(1 - \text{confidence}(n, x)\bigr) \;+\; w_3 \cdot \text{distance}\bigl(n, h(n)\bigr) \tag{PF}
\end{equation}

where $\text{severity}(x) \in [0, 1]$ is the assigned severity of exception $x$ (higher values indicate greater potential impact), $\text{distance}(n, h(n))$ is the number of hops from $n$ to the human anchor in the parent chain, and $w_1, w_2, w_3$ are provisioned constants satisfying $w_1 > w_2 > w_3 \geq 0$. The highest-priority trigger fires first, enforced by the Bailout Monitor before any state transition is committed to the Forensic Ledger.

\bigskip
\noindent\textbf{Secondary trigger: explicit \texttt{bailout\_signal}.} In addition to TC, the Bounds Engine may emit an explicit \texttt{bailout\_signal} at any time, regardless of confidence scores. This handles the case where the Bounds Engine detects a condition — for example, a logical contradiction between two inputs within $R(n)$ — that it can recognize as unresolvable even though both inputs are individually within the operating range. The \texttt{bailout\_signal} bypasses TC entirely and immediately initiates T3:

\begin{equation}
\texttt{bailout\_signal}(x) \;\implies\; \text{T3 fires for node } n \text{ with exception } x. \tag{TC-secondary}
\end{equation}

\bigskip
\noindent\textbf{The no-discretion property.} The \bounds{} layer does not exercise discretion over the trigger condition:

\begin{enumerate}
    \item The evaluation of $\text{TRIGGER}(n, x)$ is a read-only query against the operating range definition $R(n)$, not a judgment call made by a machine actor.
    \item The result of the query directly drives the state machine transition without an intervening decision step. There is no machine-actor gate between trigger evaluation and protocol activation.
    \item The confidence function is provisioned at node initialization and is immutable at runtime. No machine actor may adjust $\tau$ to suppress escalation.
\end{enumerate}

These constraints ensure that the trigger condition is architecturally enforced rather than operationally optional.

% --------------------------------------------------------------
\subsection{Escalation Algorithm}\label{sec:escalation-algorithm}
% --------------------------------------------------------------

The escalation algorithm governs the propagation of an exception from the originating node to $h(n)$ along the parent chain. It is embedded in the \fact{} layer and executed by the Bailout Monitor. The algorithm is invoked atomically upon transition T4 (BOUNDED $\rightarrow$ ESCALATING) and runs to completion or to terminal timeout.

\bigskip
\noindent\textbf{Algorithm \texttt{Escalate$(n, e)$}.}

\begin{enumerate}
    \item \textbf{Initialize.} Set $\text{current} \leftarrow n$, $\text{chain} \leftarrow [\;]$, $\text{attempt} \leftarrow 0$. Append $n$ and its diagnostic context to \text{chain}. Record \texttt{ESCALATING\_START} in the Forensic Ledger with the originating node, timestamp, and exception identifier $e$.

    \item \textbf{Select pathway.} Query the Pathway Health Registry for the currently functional pathway to $\alpha(\text{current})$. If no functional pathway exists, mark all pathways \texttt{DEGRADED} and advance to the fallback pathway selection. If all pathways are degraded, record \texttt{PATHWAY\_EXHAUSTED} and terminate with RESOLVED.

    \item \textbf{Send to parent.} Send exception $e$ with chain $\text{chain}$ to $\alpha(\text{current})$ via the selected pathway. Start timer $T_{\text{prop}}$.

    \item \textbf{Wait for acknowledgment.} Await ACK from $\alpha(\text{current})$ before $T_{\text{prop}}$ expires. If ACK is received: if $\alpha(\text{current}) = h(n)$, transition the originating node to HUMAN\_ACKNOWLEDGED; else set $\text{current} \leftarrow \alpha(\text{current})$, append $\text{current}$ to $\text{chain}$, and goto Step 2.

    \item \textbf{Timeout on parent.} If $T_{\text{prop}}$ expires without ACK: increment $\text{attempt}$. If $\text{attempt} < N_{\max}$, double $T_{\text{prop}}$ (capped at $T_{\text{cap}}$) and goto Step 2. If $\text{attempt} \geq N_{\max}$, record \texttt{PARENT\_UNREACHABLE} in the Forensic Ledger and attempt delivery via an alternate functional pathway tracked in the Pathway Health Registry. If no alternate pathway is available, terminate with RESOLVED and tag \texttt{PATHWAY\_EXHAUSTED}.

    \item \textbf{Human acknowledgment.} Upon receipt of $e_{\text{ack}}$ from $h(n)$, transition the originating node to HUMAN\_ACKNOWLEDGED. Record the full propagation chain, cumulative latency, and acknowledgment timestamp in the Forensic Ledger.

    \item \textbf{Human directive.} Await directive from $h(n)$: \texttt{RESOLVE} records the directive in the authorization ledger and transitions to RESOLVED; \texttt{REJECT} records the rejection and transitions to RESOLVED; \texttt{ACK} returns the node to IDLE. All directives are immutable and cryptographically sealed upon recording.
\end{enumerate}

\begin{figure}[htbp]
\centering
\small
\begin{tabular}{@{}ll@{}}
\toprule
Step & \multicolumn{1}{c}{Action} \\
\midrule
1 & Initialize chain, set current $\leftarrow n$ \\
2 & Select lowest-latency functional pathway \\
3 & Send exception to $\alpha$(current) with chain \\
4 & Wait for ACK within $T_{\text{prop}}$; on ACK: recurse or terminate \\
5 & On timeout: exponential backoff, retry up to $N_{\max}$ \\
6 & On $N_{\max}$ exhaust: alternate pathway or \texttt{PATHWAY\_EXHAUSTED} \\
7 & On $e_{\text{ack}}$ from $h(n)$: enter HUMAN\_ACKNOWLEDGED \\
8 & Await directive; record and transition to RESOLVED \\
\bottomrule
\end{tabular}
\caption{Escalation algorithm action table for \texttt{Escalate}$(n, e)$.}
\label{tab:escalate-actions}
\end{figure}

The algorithm maintains an append-only propagation chain $\text{chain}$ — a list of node identifiers and their diagnostic metadata — that is carried with the exception at each hop. This chain is recorded in the Forensic Ledger upon protocol termination and provides the complete attributable path from the originating node to the human anchor. It is the forensic artifact that enables post-hoc accountability tracing for any exception that reached the human principal.

\bigskip
\noindent\textbf{Pathway Health Registry.} The Pathway Health Registry (PHR) is a runtime data structure maintained by the Bailout Monitor that tracks the availability, latency, and error rate of each provisioned communication pathway to the human anchor. Before each \texttt{SelectPathway} call, the registry is queried and the pathway with the lowest estimated latency among currently functional pathways is selected. Pathways that fail $K$ consecutive health checks (default $K = 3$) are marked \texttt{DEGRADED} and excluded from pathway selection until a health check succeeds. A background health-ping thread updates the PHR at intervals of $T_{\text{health}}$ (default: 30 seconds), operating with the same scheduling priority as the Bailout Monitor.

% --------------------------------------------------------------
\subsection{Bypass Mechanism}\label{sec:bypass}
% --------------------------------------------------------------

The parent-chain bypass is the structural property that ensures the Bailout Protocol's escalation pathway contains no machine actors as terminal or intermediate absorption points between the originating node and $h(n)$. Formally:

\begin{align}
\text{BP} &: \forall n \in \text{nodes},\; \forall e \in \text{exceptions},\; \text{state}(n, e) = \text{ESCALATING} \implies \text{terminus}\bigl(\text{chain}(n, e)\bigr) = h(n). \tag{BP}
\end{align}

That is, for every node and every exception in ESCALATING state, the terminus of the propagation chain is the human anchor $h(n)$ — not any intermediate machine actor. The chain skips directly from each node to its parent, and from each parent to its parent, until $h(n)$ is reached. No node along the path receives the exception as a terminus: each node is a forwarder, not a resolver.

\bigskip
\noindent\textbf{Why bypass is architecturally necessary.} The bypass is not an optimization — it is a structural necessity that follows from the upstream accountability guarantee of the \bxthree{} Framework. Any machine actor $m$ that could absorb an exception in ESCALATING state would become a terminal node for that exception, and the accountability chain would terminate at $m$ rather than at a human. If $m$ then fails, or if $m$'s resolution is incorrect or insufficient, no human principal is present in the chain to detect or correct the failure. The system compounds failure instead of surfacing it. This is the precise failure mode the Bailout Protocol is designed to prevent.

The parent-chain bypass eliminates this by architectural constraint: there is no machine actor in the communication path, so interception, redirection, and suppression are not architecturally possible. This is analogous to a hardware interrupt line hardwired to a designated handler — just as a hardware interrupt cannot be intercepted or redirected by user-space code, the bailout escalation cannot be intercepted or redirected by any agent-level component.

The bypass is enforced by two cooperating mechanisms:

\begin{enumerate}
    \item[\textbf{Fact Layer gate.}] The \fact{} layer is the only mechanism capable of authorizing state transitions in the \bxthree{} Framework. The Bailout Monitor's pre-commit hook rejects any ledger entry for a resolution where the underlying trigger satisfies TC and $T_{\text{bounds}}$ has expired. This prevents the Bounds Engine from issuing a retroactive resolution after timeout to suppress warranted escalation. The \fact{} layer gate operates identically on all machine actors: there is no privileged machine actor with the ability to bypass this check.

    \item[\textbf{Immutable parent-pointer and human-root architecture.}] The parent pointer $\alpha(n)$ and the human-root identifier $h(n)$ are stored in \fact-layer-managed memory and are read-only at the machine-actor level. A machine actor cannot redirect the parent chain, cannot insert itself as an intermediate resolver, and cannot modify the $humanRootId$ field that terminates the chain. Any attempted modification requires \purpose{} layer authorization with a corresponding ledger entry — which cannot be completed without human approval if the modification would affect the bailout terminus.
\end{enumerate}

Together, these mechanisms ensure that BP holds for all $n \in \text{nodes}$ under all execution conditions, including partial node failure, network partition, and deliberate attempt by a compromised machine actor to redirect the bailout pathway.

Conventional escalation hierarchies — including tiered oversight models that permit absorption at intermediate tiers — make human oversight a high-tier option rather than a compulsory terminus. The Bailout Protocol forbids this absorption by architectural constraint: $h(n)$ is not one option among many — it is the \emph{only} permissible terminus for every unresolved exception.

% --------------------------------------------------------------
\subsection{Timeout Handling}\label{sec:timeout-handling}
% --------------------------------------------------------------

Timeouts in the Bailout Protocol serve a liveness function: they prevent the protocol from permanently stalling when a parent node is unreachable, a communication pathway is unresponsive, or the human anchor is temporarily unavailable. Each timeout has a defined scope, a defined scope-exit action, and an escalation path that prevents deadlock.

\bigskip
\noindent\textbf{Timeout definitions.}

\begin{enumerate}
    \item[\textbf{$T_{\text{bounds}}$ — Bounds Engine resolution timeout.}] The maximum time a node may remain in BOUNDED state attempting autonomous resolution. Default: 60 seconds, provisioned per node class according to the expected latency of the node's resolution procedure. If $T_{\text{bounds}}$ expires without a \fact-layer-approved resolution, the node transitions to BAIL\_PENDING $\rightarrow$ ESCALATING via T3 and T4. The \fact{} layer cannot extend $T_{\text{bounds}}$ unilaterally; any extension requires human authorization via a separate ledger entry recorded before $T_{\text{bounds}}$ expires.

    \item[\textbf{$T_{\text{prop}}$ — Per-hop propagation timeout.}] The maximum time node $n$ waits for an acknowledgment from its parent after sending an exception via \texttt{SendToParent}. Default initial value: 5 seconds. If $T_{\text{prop}}$ expires without ACK, the algorithm retries with exponential backoff: $T_{\text{prop}} \leftarrow \min(2 \cdot T_{\text{prop}},\; T_{\text{cap}})$. This backoff continues until either an ACK is received or $N_{\max}$ retries have been exhausted at $T_{\text{cap}}$.

    \item[\textbf{$T_{\text{cap}}$ — Maximum propagation timeout value.}] The ceiling value for the exponential backoff of $T_{\text{prop}}$. Default: 60 seconds. After $N_{\max}$ consecutive retries at $T_{\text{cap}}$ without ACK, the algorithm fires a \texttt{parent\_unreachable} timeout event, which propagates to $h(n)$ via the next available functional pathway tracked by the Pathway Health Registry.

    \item[\textbf{$T_{\text{human}}$ — Human acknowledgment timeout.}] The maximum time the human anchor may remain in HUMAN\_ACKNOWLEDGED state without issuing a directive. Default: 300 seconds (5 minutes), configurable per deployment based on the expected human response time for the node's operational domain. If $T_{\text{human}}$ expires, the protocol issues a reminder notification to $h(n)$ and transitions the originating node to RESOLVED with an \texttt{unresolved} tag. The node becomes quiescent: no autonomous retry occurs, no further actions are attempted, and the \texttt{unresolved} tag in the Forensic Ledger signals that human re-engagement is required to clear the state. The protocol does not re-escalate automatically; automatic re-escalation would create a potential denial-of-service condition against the human anchor.
\end{enumerate}

\bigskip
\noindent\textbf{Liveness guarantee.} Together, $T_{\text{bounds}}$, $T_{\text{prop}}$, $T_{\text{cap}}$, and $T_{\text{human}}$ define a finite upper bound on the wall-clock time from trigger firing (T1) to human acknowledgment under any combination of parent unavailability and pathway degradation:

\begin{equation}
T_{\max} \;=\; T_{\text{bounds}} \;+\; N_{\max} \cdot T_{\text{cap}} \;+\; T_{\text{human}} \tag{T-max}
\end{equation}

A deployment that cannot guarantee human acknowledgment within the required $T_{\max}$ under all failure scenarios must provision additional or higher-bandwidth pathways, or must reduce $T_{\text{prop}}$ and $T_{\text{cap}}$ to bring the worst-case bound within the required limit. This analysis is a required part of deployment certification for any \bxthree{} system operating in a regulated or safety-critical environment.

\bigskip
\noindent\textbf{Timeout interaction with the bypass mechanism.} The bypass mechanism ensures that timeouts cannot be weaponized by a machine actor to delay or suppress escalation. Specifically: (a) $T_{\text{bounds}}$ is enforced by the \fact{} layer pre-commit hook, not by the \bounds{} layer; a machine actor cannot artificially extend $T_{\text{bounds}}$ to prevent the transition to ESCALATING. (b) $T_{\text{prop}}$ backoff is calculated by the Bailout Monitor, not by any machine actor; the backoff curve is defined in the protocol specification and cannot be altered at runtime. (c) If all pathways timeout at $T_{\text{cap}}$, the protocol resolves with \texttt{PATHWAY\_EXHAUSTED} — this is a terminal RESOLVED state, not a suppression of escalation. The \texttt{PATHWAY\_EXHAUSTED} tag in the Forensic Ledger ensures that the failure of the pathway is recorded as an incident requiring human remediation, not as a silently suppressed exception.
